\documentclass[11pt,a4paper]{article}
\usepackage[T1]{fontenc}
\usepackage[utf8]{inputenc}
\usepackage{lmodern}
\usepackage{amsmath,amssymb}
\usepackage[margin=25mm]{geometry}
\usepackage[hidelinks]{hyperref}
\hypersetup{pdftitle={One action constant from composition and spatial rotations},pdfauthor={navstokgap research working notes}}
\newcommand{\tightlist}{\setlength{\itemsep}{0pt}\setlength{\parskip}{0pt}}
\setlength{\emergencystretch}{3em}
\setcounter{secnumdepth}{0}
\begin{document}
\hypertarget{one-action-constant-from-composition-and-spatial-rotations}{%
\section{One action constant from composition and spatial
rotations}\label{one-action-constant-from-composition-and-spatial-rotations}}

\textbf{After refereeing (Fable, 2026-09-27).} The identity
\([L_x,L_y]=i\sum_ih_iL_{iz}\) is exact, so Theorem 1's closure
requirement restates \(h_i=h\); its content is the premise it names,
that the measured additive angular momentum is the same multiple of the
rotation generator for every body. The genuine dynamical result is
equation (10): if each body's Newton equations hold with its own
constant in a shared product, any pair with a nonzero mixed force
derivative \(\partial_{ia}\partial_{jb}V\) must have \(h_i=h_j\).
Interaction forces a common constant, with no appeal to centrality or
angular momentum; \(h=0\) stays admissible. Theorem 2 is a special case,
with the sign of (8) corrected.

\textbf{Result, 2026-09-27 (round 2).} Independent constituent canonical
algebras with constants \(h_i\) obey \([L_x,L_y]=i\sum_i h_iL_{iz}\).
Requiring the \textbf{mechanical sum}
\(\mathbf L=\sum_i\mathbf q_i\times\mathbf p_i\) to obey
\([L_x,L_y]=ihL_z\) as an identity on the full constituent algebra
forces \(h_i=h\) for every constituent (Theorem 1). Equivalently, at
\(h\ne0\), requiring \(\mathbf L/h\) to generate the same spatial
rotation on every constituent fixes the same normalization. Admissible
pair composition propagates this equality through a connected graph.
Central interactions provide a second route if conservation of the
mechanical sum is imposed under a common commutator evolution (Theorem
2).

Both conclusions admit a common zero value. Spatial rotation covariance
alone permits unequal \(h_i\): for nonzero constants the dimensionless
generator is \(\sum_i\mathbf L_i/h_i\). The additional physical premise
identifies this generator with the measured additive angular momentum
divided by one action constant. Integer and half-integer representation
labels constrain angular momenta relative to that constant; selecting
its positive magnitude requires further physical input.

These are elementary conditional results, with written proofs below.
They test the user's composition proposal following
\href{rivero-1998-conjecture-central-forces.md}{round 1, §§4--5}
(full-read). Here \textbf{\(h\) is the reduced phase constant}, so
identification with quantum mechanics would give \(h=\hbar\) and
\(h_P=2\pi h\). The algebra permits signed real constants; positive
floors use \(|h|\). All momenta are measured quantities of motion, with
fixed mass and unit conventions throughout.

\hypertarget{the-constituent-algebra-and-what-rotation-must-mean}{%
\subsection{1. The constituent algebra and what rotation must
mean}\label{the-constituent-algebra-and-what-rotation-must-mean}}

Take finitely many labelled bodies, \(i=1,\ldots,N\), with masses
\(m_i>0\), three spatial coordinates and three momenta. Assume an
associative unital complex algebra with self-adjoint generators
satisfying

\[
[q_{ia},q_{jb}]=[p_{ia},p_{jb}]=0,\qquad
[q_{ia},p_{jb}]=i\delta_{ij}\delta_{ab}h_i1,
\quad a,b\in\{x,y,z\}. \tag{1}
\]

The scalars \(h_i\in\mathbb R\) remain fixed when the body is put in a
composite. Distinct constituent algebras commute. Assume also that the
normally ordered monomials \(q^\alpha p^\beta\) are linearly
independent: we retain the full tensor product of canonical polynomial
algebras, including its commutative factors when some \(h_i=0\). This is
the precise independence used in coefficient comparisons. Statistical
independence of states is unnecessary; entangled states are allowed. For
unbounded operator realizations all displayed identities are on a common
invariant dense domain representing these polynomials faithfully.

Define, with \(\epsilon_{xyz}=1\),

\[
L_{ia}=\sum_{b,c}\epsilon_{abc}q_{ib}p_{ic},\qquad
L_a=\sum_iL_{ia}. \tag{2}
\]

Each term in (2) has unambiguous ordering because \(q_{ib}\) and
\(p_{ic}\) commute for \(b\ne c\). The additivity and mechanical
normalization in (2) are hypotheses. Independently rescaling each body's
momentum would change its measured quantity of motion and these
hypotheses.

Three distinct requirements can now be stated precisely.

\begin{itemize}
\tightlist
\item
  \textbf{Spatial covariance:} the map \(q_i\mapsto Rq_i\),
  \(p_i\mapsto Rp_i\) is an algebra automorphism for every
  \(R\in SO(3)\).
\item
  \textbf{Mechanical closure:} for one real \(h\),
  \([L_a,L_b]=ih\sum_c\epsilon_{abc}L_c\) on the full algebra.
\item
  \textbf{Mechanical generator identification:} for \(h\ne0\), the
  generator \(K_a=L_a/h\) satisfies
  \([K_a,q_{ib}]=i\sum_c\epsilon_{abc}q_{ic}\) and
  \([K_a,p_{ib}]=i\sum_c\epsilon_{abc}p_{ic}\) for every body.
\end{itemize}

The last convention gives \(U(\theta)=\exp(-i\theta K_a)\) and
\(U(\theta)^\dagger q_{ib}U(\theta) =q_{ib}-\theta\sum_c\epsilon_{abc}q_{ic}+O(\theta^2)\),
the usual active rotation of the vector's components about axis \(a\).
Integration to a group representation additionally needs a strongly
continuous unitary action with these generators. Lie commutator
identities alone leave those domain and integrability obligations open.
For \(h\ne0\), such an action is a representation of \(SU(2)\); descent
to \(SO(3)\) additionally requires a full \(2\pi\) rotation to act as
the identity.

The round-1 phase \(e^{iL\theta/h_i}\) yields (1) only after an
additional canonical observable/quantization prescription. A single
orbit's angular descent condition by itself supplies neither the tensor
algebra nor its mechanical generator identification. The present theorem
begins at that explicit algebraic extension.

\hypertarget{theorem-1-mechanical-closure-forces-a-shared-constant}{%
\subsection{2. Theorem 1: mechanical closure forces a shared
constant}\label{theorem-1-mechanical-closure-forces-a-shared-constant}}

\textbf{Theorem 1.} Under (1)--(2) and the monomial independence
hypothesis, mechanical closure holds for a scalar \(h\) if and only if
\(h_i=h\) for all \(i\). It suffices to impose the single identity
\([L_x,L_y]=ihL_z\). For \(h\ne0\), mechanical generator identification
is also equivalent to \(h_i=h\) for all \(i\).

\emph{Proof.} The product rule for commutators gives

\[
[L_{ia},q_{jb}]=i\delta_{ij}h_i\sum_c\epsilon_{abc}q_{ic},\qquad
[L_{ia},p_{jb}]=i\delta_{ij}h_i\sum_c\epsilon_{abc}p_{ic}. \tag{3}
\]

For example, writing \(L_{ix}=q_{iy}p_{iz}-q_{iz}p_{iy}\) and
\(L_{iy}=q_{iz}p_{ix}-q_{ix}p_{iz}\), their commutator has exactly the
surviving terms

\[
-i h_i q_{iy}p_{ix}+i h_i q_{ix}p_{iy}=ih_iL_{iz}.
\]

Different bodies commute, so closure of the sum is equivalent to

\[
\sum_i(h_i-h)L_{iz}=0. \tag{4}
\]

The monomial \(q_{ix}p_{iy}\) in (4) occurs only in \(L_{iz}\), with
coefficient \(h_i-h\). Independence makes that coefficient zero for each
\(i\), including a constituent with \(h_i=0\). Conversely, equal
constants give closure in all three directions by the same calculation.

For generator identification, (3) gives, for example,
\([L_x,q_{iy}]=ih_iq_{iz}\), whereas the required action is
\(ihq_{iz}\). Independence makes \(q_{iz}\ne0\), hence \(h_i=h\). The
converse follows directly from (3). \(\square\)

A restricted observable space needs the corresponding separation
hypothesis in place of monomial independence: each relation
\(\sum_i c_iL_{ia}=0\) must force every \(c_i=0\). Compression to a
subspace where all orbital \(L_i\) vanish makes closure vacuous for
arbitrary \(h_i\). A single orbit, expectation value or rigidly
constrained configuration can likewise lose the necessary separation.
The theorem uses observable identities valid for all admitted
preparations and configurations, with the constituents still
individually accessible.

\textbf{Transitivity.} Form a graph whose vertices are bodies or species
with persistent constants \(h_i\). Put an edge between two vertices when
their full pair algebra is admitted and its mechanical sum is required
to have one closure constant. Theorem 1 gives equality on each edge and
therefore throughout each connected component. Connectedness gives one
constant for the whole class, without assuming that every pair interacts
directly. Disconnected components remain unrelated unless a common
composite or generator condition links them. A context-dependent \(h_i\)
would require a further preparation-equivalence premise before this
argument could compare contexts.

\textbf{Role of three dimensions.} A single planar rotation generator
has \([L_z,L_z]=0\), which tests no normalization. In three dimensions
the non-abelian bracket in (4) compares the mechanical sums internally.
Requiring the correct action on each \(q_i,p_i\) already compares
normalizations in two dimensions by (3). Three dimensions strengthen the
closure test; generator identification itself needs only one nontrivial
rotation plane.

\hypertarget{rotation-covariance-and-the-interaction-test}{%
\subsection{3. Rotation covariance and the interaction
test}\label{rotation-covariance-and-the-interaction-test}}

\hypertarget{covariance-with-unequal-constants}{%
\subsubsection{Covariance with unequal
constants}\label{covariance-with-unequal-constants}}

For arbitrary \(h_i\), each scalar canonical block in (1) is rotation
invariant. An explicit associative realization on polynomials is

\[
f*g=\mu\exp\left[\frac i2\sum_i h_i P_i\right] (f\otimes g),\qquad
P_i=\sum_a\left(\partial_{q_{ia}}\otimes\partial_{p_{ia}}
-\partial_{p_{ia}}\otimes\partial_{q_{ia}}\right). \tag{5}
\]

The exponential terminates on polynomials. The constant derivative
operators commute, so either association of three factors gives the same
exponential of the sum over their three pairings. Each \(P_i\) is
invariant under simultaneous spatial rotations. Thus (5) supplies
covariance for every tuple \((h_1,\ldots,h_N)\).

If all \(h_i\ne0\), define

\[
K_a=\sum_i L_{ia}/h_i.
\]

Equations (3) give \([K_a,K_b]=i\epsilon_{abc}K_c\) (summed \(c\)) and
the correct rotations of all constituent observables. Concretely, on
\(L^2(\mathbb R^{3N})\), \(p_i=-ih_i\nabla_i\) yields
\(\mathbf K=-i\sum_i\mathbf q_i\times\nabla_i\), which integrates to
simultaneous scalar rotations for arbitrary nonzero \(h_i\). The
identification \(K_a=L_a/h\) adds exactly the normalization tested in
Theorem 1. With some zero factors, spatial covariance still holds as an
automorphism; commutators with their central observables vanish, so
their rotations require an outer derivation.

\hypertarget{central-interactions-and-conservation}{%
\subsubsection{Central interactions and
conservation}\label{central-interactions-and-conservation}}

For classical Newton dynamics about a fixed inertial origin,
\(\dot q_i=p_i/m_i\) and \(\dot p_i=\sum_{j\ne i}F_{ij}\), the strong
central form of the Third Law is \(F_{ji}=-F_{ij}\) with \(F_{ij}\)
parallel to \(q_i-q_j\). It gives

\[
\dot{\mathbf L}
=\sum_{i<j}(q_i-q_j)\times F_{ij}=0. \tag{6}
\]

Equal and opposite forces together with centrality do the work here.
Equation (6) uses neither a commutator nor an action constant.
Transferring this conservation law to the deformed observable algebra
needs a specified evolution law.

\textbf{Theorem 2 (conservation propagates equality along interacting
pairs).} Assume (1), faithful configuration observables on open sets,
and a Hamiltonian

\[
H=\sum_i\frac{p_i^2}{2m_i}+\sum_{i<j}V_{ij}(r_{ij}),\qquad
r_{ij}=|q_i-q_j|. \tag{7}
\]

Assume commutator differentiation \([p_{ia},V]=-ih_i\partial_{q_{ia}}V\)
on a common domain. This holds algebraically for polynomial potentials,
or on a smooth domain away from collisions for the coordinate
realization. Give time evolution one fixed nonzero action normalization
\(h_t\), \(\dot A=[A,H]/(ih_t)\). Require conservation of the mechanical
sum for every admitted pair interaction, with other pair couplings
independently switchable off. For each tested edge require
\(V_{ij}'(r_{ij})\ne0\) on some open set allowing \(q_i\times q_j\ne0\).
Then conservation forces \(h_i=h_j\) on that edge, and therefore one
value per connected interaction component.

\emph{Proof.} The kinetic energies commute with each \(L_i\). For one
pair, \(\nabla_j V_{ij}=-\nabla_i V_{ij}\), and centrality gives
\((q_i-q_j)\times\nabla_i V_{ij}=0\). Consequently

\[
[\mathbf L,V_{ij}]
=i(h_i-h_j)q_i\times\nabla_i V_{ij}
=+i(h_i-h_j)\frac{V_{ij}'(r_{ij})}{r_{ij}}q_i\times q_j. \tag{8}
\]

The configuration factor is a nonzero observable on the stated open set.
Conservation, equivalent to \([\mathbf L,H]=0\), forces \(h_i=h_j\).
When constants agree on every interacting edge, (8) proves the converse
conservation statement as well. \(\square\)

One central harmonic interaction suffices as a polynomial test:
\(V_{ij}=k_{ij}|q_i-q_j|^2/2\), \(k_{ij}>0\), gives the exact defect
\(+ik_{ij}(h_i-h_j)q_i\times q_j\) (sign corrected 2026-09-27 after a
Fable review; the conclusion is unaffected). This permits full
configuration variation and keeps every step inside the polynomial
algebra. Restriction to collinear configurations or to a pair's fixed
centre-of-mass frame would remove this particular torque test;
conservation about arbitrary fixed origins restores its stated scope. An
external central potential for each isolated body gives no such exchange
test and allows different constants.

Theorem 2 compares the \(h_i\) but leaves their common value relative to
\(h_t\) undetermined. Demanding the full Newton equations with the
physical Hamiltonian (7) gives the stronger identities

\[
\dot q_i=\frac{h_i}{h_t}\frac{p_i}{m_i},\qquad
\dot p_i=-\frac{h_i}{h_t}\nabla_i V. \tag{9}
\]

Requiring \(\dot q_i=p_i/m_i\) for unrestricted momenta forces
\(h_i=h_t\). This route already assumes a nonzero inner-commutator
clock. Alternatively, assigning each body its own denominator and
retaining all Newton equations must respect the shared product: for any
derivation \(D\) with \(Dp_{ia}=-\partial_{ia}V\),

\[
0=D[p_{ia},p_{jb}]
=i(h_i-h_j)\partial_{ia}\partial_{jb}V,\qquad i\ne j. \tag{10}
\]

A nonzero mixed force derivative again equates the constants. This is
the main dynamical result of the note (reviewer's assessment): it needs
no centrality and no angular momentum, and it holds for every
interacting pair. Equation (10) follows by applying the Leibniz rule to
the zero cross commutator; it exposes the compatibility obligation of
separate denominators. The common zero branch satisfies (10) and
supports ordinary classical Hamiltonian evolution by a Poisson
derivation.

\hypertarget{what-happens-at-zero-and-what-spin-can-add}{%
\subsection{4. What happens at zero, and what spin can
add}\label{what-happens-at-zero-and-what-spin-can-add}}

\textbf{Corollary (universality and positivity).} Theorem 1 forces a
common real value, with \(h=0\) allowed. If one constituent has an
independently established \(h_i\ne0\), connected composition propagates
its nonzero value to the entire component. The positivity input lies in
that reference constituent. Theorem 2 likewise permits all \(h_i=0\);
its commutator clock then gives trivial evolution. Classical nontrivial
dynamics uses the Poisson bracket instead.

For the common family (5), all real \(h\) give associative algebras. At
\(h=0\) the product is ordinary multiplication. Spatial rotations remain
non-abelian: their derivations are \(D_a f=\{f,L_a\}_{\rm P}\), and
classical angular momenta obey
\(\{L_a,L_b\}_{\rm P}=\epsilon_{abc}L_c\). The Poisson algebra supplies
the rotation structure while the multiplication algebra is commutative.
For polynomial observables the common Moyal family has

\[
\lim_{h\to0}\frac{[f,g]_{*_h}}{ih}=\{f,g\}_{\rm P}. \tag{11}
\]

Requiring a nontrivial rotation to be implemented specifically by
\(\exp(-i\theta L_a/h)\) selects \(h\ne0\) in the premise. Its formula
requires division by \(h\); classical rotations remain perfectly defined
by their derivations at zero. A strongly continuous unitary rotation
representation alone also leaves this choice open: pullback rotates
classical phase-space functions in \(L^2\) using differential
generators, while measured classical angular momenta act by
multiplication.

For \(h\ne0\), supply self-adjoint angular generators integrating to a
unitary \(SU(2)\) representation. Its irreducible sectors have

\[
J^2=h^2j(j+1),\qquad J_z=hm,\qquad
j\in\{0,\tfrac12,1,\tfrac32,\ldots\},\quad
m=-j,-j+1,\ldots,j. \tag{12}
\]

Here \(J\) denotes total angular momentum, including any supplied
internal spin. To recall the dimensionless argument, put \(K=J/h\) and
\(K_\pm=K_x\pm iK_y\). Then \([K_z,K_\pm]=\pm K_\pm\) and
\(K_\mp K_\pm=K^2-K_z^2\mp K_z\). In an irreducible unitary sector,
positivity bounds the ladder above and below. Its highest weight \(j\)
has \(K^2=j(j+1)\); the lowest weight is \(-j\), so \(2j\) is a
nonnegative integer. A full rotation acts by \(e^{-2\pi im}=(-1)^{2j}\);
\(SO(3)\) descent retains integer \(j\).

For ordinary scalar orbital wavefunctions, \(L=q\times(-ih\nabla)\) has
integer orbital labels and spherical harmonics. The canonical angular
algebra and this orbital spectrum are standard; see
\href{https://www.damtp.cam.ac.uk/user/tong/qm/qmhtml/S4.html}{Tong,
\emph{Quantum Mechanics}, §§4.1.1--4.1.2} (passage, commutator and
spherical-harmonic derivations read). Half-integer spin requires
additional internal degrees of freedom or an explicitly different state
space. The central-force orbital construction supplies no such degrees
of freedom by itself. If independent spin observables commute with the
orbital variables, requiring \(J=L+S\) to rotate those variables still
gives (3), hence the same orbital universality result. Requiring it also
to rotate each nontrivial spin algebra fixes the spin normalization in
the same way.

Equation (12) supplies dimensionless spectral ratios. The Casimir is
discrete; its factor \(j(j+1)\) is integer for integer \(j\) and can be
fractional for half-integer \(j\). Conditional on fixed \(|h|>0\), the
smallest nonzero \(J^2\) among all integer sectors is \(2h^2\), and
among all \(SU(2)\) sectors it is \(3h^2/4\). A chosen representation
can omit these sectors; the scalar sector \(j=0\) always remains an
allowed representation. Both spectral factors leave the dimensional
value of \(h\) free.

More explicitly, rescaling \(J\mapsto\lambda J\), \(h\mapsto\lambda h\)
for \(\lambda>0\) preserves the generators \(K\), all rotation matrices,
and the integer/half-integer labels. A measured nonzero magnitude
\(\ell\) and a known \(j>0\) would fix \(|h|=\ell/\sqrt{j(j+1)}\), using
that dimensional measurement as input. With only \(\ell\) fixed,
unbounded \(j\) allows values tending to zero. Several specified
component and Casimir eigenvalues can further constrain \(j\) and their
ratios; this constitutes additional spectral data. The scale freedom of
the representation conditions alone remains.

\hypertarget{comparison-with-the-two-composition-and-dimensional-results}{%
\subsection{5. Comparison with the two composition and dimensional
results}\label{comparison-with-the-two-composition-and-dimensional-results}}

The \href{composition-universality.md}{classical composition note,
§§1--3} (full-read) concerns a fluctuation coefficient \(\kappa(m)\ge0\)
and independent displacement variances. Its whole-body premise is

\[
\kappa(m_1+m_2)
=\frac{m_1\kappa(m_1)+m_2\kappa(m_2)}{m_1+m_2}.
\]

For all positive masses, \(f(m)=m\kappa(m)\) is additive and
nonnegative, hence linear; \(\kappa(m)=K\ge0\). A reference constituent
then supplies \(K=m_0u_0^2/\lambda_0>0\) under the stated nonzero-speed
and finite-rate premises. The rotation theorem replaces this
mass-additivity premise with mechanical generator identification and
separation of constituent observables. A connected species graph
suffices, and arbitrary positive masses need not be admitted. Both
arguments propagate a coefficient; each obtains positivity only with
additional input. Identifying \(K\) with the phase constant \(h\)
remains another physical obligation.

There is an exact parallel in centre and relative variables. Set
\(M=m_1+m_2\), \(\mu=m_1m_2/M\),

\[
R=\frac{m_1q_1+m_2q_2}{M},\quad r=q_1-q_2,\quad
P=p_1+p_2,\quad p=\frac{m_2p_1-m_1p_2}{M}.
\]

Bilinearity in (1) gives

\[
\begin{aligned}
{}[R_a,P_b]&=i\delta_{ab}\frac{m_1h_1+m_2h_2}{M},&
[r_a,p_b]&=i\delta_{ab}\frac{m_2h_1+m_1h_2}{M},\\
[R_a,p_b]&=i\delta_{ab}\frac{\mu}{M}(h_1-h_2),&
[r_a,P_b]&=i\delta_{ab}(h_1-h_2).
\end{aligned} \tag{13}
\]

Thus canonical independence of the centre and relative algebras forces
\(h_1=h_2\). The classical note has the corresponding covariance
\(\operatorname{Cov}(\Delta R,\Delta r)=\Delta(\kappa_1-\kappa_2)/M\);
for its Gaussian preparations vanishing covariance gives statistical
independence. Equation (13) concerns commuting observable algebras,
which can carry correlated or entangled states. The two independence
premises have distinct meanings despite their matching coefficients.

The \href{principia-fifth-postulate.md}{fifth-postulate note}
(full-read) places these results as follows. Its Theorem B obtains a
single Moyal constant under full affine symplectic covariance. Applying
that premise to the combined phase space already includes
transformations mixing bodies, including (13), whereas separate one-body
covariance allows one constant per factor as in (5). Theorem 1 isolates
a narrower comparison using physical angular generators; it leaves
general star-product uniqueness to the stronger covariance theorem.
Theorem A retains the common zero and nonzero families under its Newton
equations. Theorems C and E give conditional disturbance and
concentration bounds once \(|h|>0\) is supplied. Theorem D records
action rescaling. The Gaussian state restriction of B' supplies its own
\(\zeta\) in a commutative algebra; spatial rotations and the present
commutator test leave a relation \(2\zeta=|h|\) to an additional
premise.

Finally, \href{action-unit-dimensional-selection.md}{the dimensional
note, Theorem A} (full-read) applies to a universal floor \(g(c)\)
determined by fixed constants \(c_r\). If an action product exists, its
form is

\[
g(c)=\left(\prod_r c_r^{a_r}\right)F(\pi_1,\ldots,\pi_s),\qquad
\sum_r a_r d_r=(1,2,-1). \tag{14}
\]

It reduces to a pure number times the product when the fixed constants
have no independent dimensionless products. Rotation angles, group
structure constants, \(j\) and \(m\) are dimensionless and add no action
unit to (14). When the fixed constants admit no action-dimensional
product, Theorem A allows only a zero or infinite universal floor;
dimensionless rotations preserve that conclusion. Treating \(h\) as a
new independent fixed constant makes \(|h|\) the supplied unit. Deriving
a universal \(|h|\) from pre-existing constants instead requires
dimensional homogeneity and the same covariance argument as (14);
angular universality provides no value for its dimensionless factor or
proof of its positivity. For the fixed \(k_e=e^2/(4\pi\epsilon_0)\),
\(c\) and \(G\) of that note, this would read \(|h|=Ck_e/c\). The
rotation theorem leaves \(C\) undetermined. A coefficient and the
infimum of a specified action observable also need a proved relation
before (14) can be used as a floor statement.

\hypertarget{consequence-for-state}{%
\subsection{6. Consequence for STATE}\label{consequence-for-state}}

Item 2 gains conditional universality from mechanical angular-momentum
composition: one constant per connected class follows from faithful
rotation closure, or from conservation under the stated interacting
commutator dynamics. The precise additional premise beyond round 1 is
the identification of additive measured angular momentum with the common
rotation generator. The remaining necessity task is a physical premise
selecting a nonzero branch and an action unit; spin topology and Casimir
discreteness leave that obligation open. All conclusions above are
proved under their displayed hypotheses.
\end{document}
